\subsection{Repeated-Split Coverage and Budget Utilisation}
\label{sec:coverage}

A single calibration/test split cannot estimate how often the certificate is violated. We therefore repeat the patient-level split ten times on each site, calibrate on half of the eligible subjects and evaluate on the other half, and record the realised mean rate, the fraction of splits above the requested budget, the no-alarm fallback rate and the patient-level violation fraction. The supplement plots budget utilisation; Table~\ref{tab:coverage} reports the 1-per-day row.



\begin{table}[tb]
\centering
\caption{Repeated-split coverage at the 1-per-patient-day budget, ten stratified splits per site. ``Rate'' is the mean realised alarms per patient-day, ``Util.'' is rate divided by 1.0, and ``Patient viol.'' is the mean fraction of test patients above the per-patient budget. Blocked, LTT and randomised selectors had no split above the mean budget; CRC crossed it in one of ten A splits at 0.25 per day, consistent with its expected-risk rather than PAC guarantee. Blocked is the empirical-Bernstein selector; LTT is Learn-then-Test with Holm correction; Random. is the randomised adjacent-threshold mixture.}
\label{tab:coverage}
\small
\resizebox{\columnwidth}{!}{%
\begin{tabular}{lcccccc}
\toprule
 & \multicolumn{3}{c}{Site A} & \multicolumn{3}{c}{Site B} \\
\cmidrule(lr){2-4}\cmidrule(lr){5-7}
Procedure & Rate & Util. & Patient viol. & Rate & Util. & Patient viol. \\
\midrule
Blocked & $0.634$ & $0.63$ & $0.210$ & $0.630$ & $0.63$ & $0.215$ \\
CRC & $0.931$ & $0.93$ & $0.296$ & $0.935$ & $0.93$ & $0.308$ \\
Timestep & $0.931$ & $0.93$ & $0.296$ & $0.935$ & $0.93$ & $0.308$ \\
Plugin & $0.931$ & $0.93$ & $0.296$ & $0.935$ & $0.93$ & $0.308$ \\
LTT & $0.872$ & $0.87$ & $0.279$ & $0.935$ & $0.93$ & $0.308$ \\
Random. & $0.878$ & $0.88$ & $0.284$ & $0.896$ & $0.90$ & $0.302$ \\
\bottomrule
\end{tabular}}
\end{table}

Three findings follow. First, the blocked, LTT and randomised certificates exceeded the mean budget in none of the ten splits on either site; the plug-in rule and CRC each crossed it in one of ten A splits at the 0.25-per-day budget, where the plug-in empirical rate was $0.221$ against a request of $0.25$ on average. That single CRC exceedance is the expected-risk signature discussed in \cref{sec:crc}, not a failure of the blocked certificate. Second, the timestep shortcut also stayed below the mean budget in all splits. The data therefore do not support the stronger claim that the shortcut always violates the mean-rate constraint; its deficiency is that it has no patient-level guarantee, and at 1 per day its patient-violation fraction is $0.296$ versus $0.210$ for the blocked certificate. It behaves like the plug-in rule because its effective sample size is large, so it selects the most aggressive empirical threshold. Third, budget utilisation is the main engineering cost of the certificate. At 1 per day the blocked selector uses 63\% of the budget on both sites, whereas CRC uses 93\% and raises sensitivity from $0.613$ to $0.706$ at A and from $0.514$ to $0.619$ at B. Across the four budgets the blocked selector uses 63--75\% at A and 63--88\% at B; CRC uses 82--93\% at A and 93\% at B. The empirical-Bernstein LTT and the randomised mixture occupy an intermediate band, at the price of a higher patient-level violation fraction. The plug-in rule reaches 88--93\% on site A at the three tighter budgets and 93\% at 0.25 on site B, but only 70\% at 0.5 and 88\% at 2 on site B; it has no finite-sample guarantee and exceeds the budget in one of the ten A splits at 0.25. This is the trade that the original single-split table could not show: the certificate is conservative, and the conservatism can be reduced by a sharper selector or by randomisation, not by pretending that the point estimate is a certificate.

\subsection{Operational Alarm Metrics}
\label{sec:operational}

Sensitivity and mean rate do not specify how many alarms a clinician sees for each detected event. We therefore report positive predictive value (PPV), the number needed to alert (NNA), and alarms per detected event, all computed from the same replay as the policy outcomes. On the primary A split at 1 per day, the blocked logistic policy emits $0.631$ alarms per patient-day with PPV $0.060$ and NNA $16.7$; the plug-in policy emits $0.905$ alarms per day with PPV $0.050$ and NNA $20.0$. On the external B evaluation the corresponding blocked values are PPV $0.037$ and NNA $27.2$, and the plug-in values are PPV $0.032$ and NNA $31.4$. At the tighter 0.25-per-day budget the blocked policy has PPV $0.099$ and NNA $10.1$ at A, because the same number of true alarms is divided by far fewer false alarms. The absolute PPV is low because the event prevalence is low and a patient can generate many hourly alarms; the useful engineering statement is the monotone trade: a tighter budget raises PPV and lowers NNA, while a looser budget buys sensitivity with more false alarms per detected event. These metrics also explain why the same sensitivity can correspond to very different clinical workloads across sites.

\subsection{Target Adaptation in Brief}
\label{sec:adaptation-summary}

A target site that can observe unlabelled traces can match its requested rate without outcome labels. On a disjoint warm-up half of B, the most aggressive threshold whose empirical rate is at most the budget is deployed on the held-out half. At 1 per day this raises sensitivity from $0.414$ to $0.620$ and the realised rate from $0.356$ to $0.950$; at 0.5 it reaches $0.414$ at $0.356$, coinciding with the target-calibrated oracle on this grid. The rule is transductive and controls the rate rather than the risk, so it is a deployment heuristic rather than a replacement for a labelled certificate. The supplement reports the full adaptation and ablation tables.

\subsection{Ward-Level Alarm Load}
\label{sec:ward}

Per-patient rates are difficult to translate into staffing. We replay the calibrated traces on a synthetic 20-bed ward with Poisson arrivals, finite beds and a fixed hourly alarm-handling capacity. At five admissions per day and a capacity of one alarm per hour, the blocked policy exceeds capacity in $0.2\%$ of hours at 0.25 alarms per patient-day, $3.8\%$ at 1 per day and $12.7\%$ at 2 per day. At a capacity of two alarms per hour the corresponding fractions are $0.0\%$, $0.2\%$ and $1.6\%$ (see the supplementary operational figure). The peak hourly demand rises from two to three alarms, and the 95th percentile from one to two. At ten admissions per day the ward reaches $84\%$ bed occupancy and refuses $13.6\%$ of admissions at the 2-per-day operating point, so the load estimate is conditional on the arrival and capacity assumptions rather than a property of the policy alone. The simulation is a translation device, not a validation: it shows that a budget difference of a factor of two is operationally visible at the nursing-station level, and that a ward with one alarm per hour of spare capacity is stressed even at 1 alarm per patient-day.


A continuous-time queue provides a human-factors check on the hourly capacity abstraction. Using the 5-minute median response reported by Bridi et al.\ and a 15-minute mean that includes travel and confirmation \citep{bridi2014reaction}, one nurse handles the blocked 1-per-day policy with fewer than 2\% of alarms waiting more than 10 minutes. At 2 per day the same nurse leaves 9.4\% of alarms waiting more than 10 minutes, with a 95th-percentile wait of 26.8 minutes; a 5-minute mean reduces the over-10-minute fraction to 0.5\%, and a second nurse removes essentially all waiting. The full queue table is in the supplement. These delays are lower than the 60\% unanswered rate in the observational study because the queue assumes perfect work-conserving service; it is a lower bound on human delay, not a validation of the ward model. Even so, the difference between a 1- and 2-per-day budget is visible in response latency at average nurse utilisation below 20\%, because alarm bursts matter more than the daily mean.
